\begin{figure*}[t]
\centering
\includegraphics[width=\textwidth]{results/build_ablation.pdf}
\caption{Construction-control runtime relative to full (262144 cache entries), cold/warm: aggregate optimizer-time ratios and per-design ratios. Colors encode the logarithm of each ratio. $\dagger$ marks at least one TO; equal elapsed budgets do not imply equal completed work. Selected-output area effects and timeout counts appear in \Cref{tab:ablation}.}\label{fig:ablation}
\end{figure*}

\begin{figure}[t]
\centering
\includegraphics[width=.90\linewidth]{results/round_progress.pdf}
\caption{Selected-output reductions from $N_e$ at caps 1 and 6. Cold/warm curves separate when different; intermediate caps were not measured.}\label{fig:rounds}
\end{figure}
\section{Experimental Results}

\MfsComparisonTable
\MainResultTable
\subsection{Setup}\label{sec:setup}
We evaluate sixteen designs with the ASAP7 genlib~\cite{asap7}, reporting area in $\mu m^2$ and worst-output delay in ps. We use C++/mockturtle~\cite{EPFLLibraries}, Rust/\textit{egg}~\cite{willsey2021egg}, and ABC~\cite{ABCTool}.

All experiments fix the initial delay bound across rounds (\Cref{eq:objective}), with whole-window synthesis and unmap/remap exploration disabled.

For \texttt{mfs3}, both tools use the same 40-entry subset with at most six inputs per cell, and byte-identical area-oriented \texttt{emap} inputs. The evaluated \texttt{mfs3} implementation restricts cell truth tables and node simulation to six inputs. Preprocessing experiments and ablations use all 47 entries for initial mapping and optimization.

The direct input $N_i$ maps the raw AIG without ABC optimization. For $N_e$, the flow repeats \texttt{fraig; resyn2rs; fraig; compress2rs} until the AND count stops decreasing, with a 20-iteration cap, then applies area-oriented \texttt{emap}. Let $A_e=A(N_e)$ and $D_e=D(N_e)$. The database does not alter initial mapping; benchmarks and preprocessing remain unchanged.

Cuts use $K=6$, $B=12$, and the L4 area threshold $\theta=0$. Care uses an eight-input reconvergence window and a cumulative 60\,s budget, then full care; SAT refinement is disabled. The three-level mapped database has area cap $0.50$. SOP/DSD starts at four leaves; Shannon and both Davio forms start at three. Each native structure cache holds at most 262144 entries; pre-mapping reuse is keyed by the exact AIG state. Temporary mapped networks share an immutable cell library. All statistical deferrals are rechecked; only strict area dominance suppresses live pricing. Independent output-polarity factoring and exact AIG synthesis are disabled. Equality saturation allows two iterations and 500 e-nodes; extraction permits eight realizations per class, 20000 combinations and 4000 visits. Chain composition uses two-step blocks, at least six steps and cross-root sharing. Both round caps allow early stopping.

Experiments run on two Xeon Gold 6226R sockets (32 cores), Ubuntu 22.04.5, GCC 11.2.0, Rust 1.99.0 and ABC 1.01. A shared host runs up to sixteen single-worker ePMapping optimizations and eight independent ABC jobs. Both tools have a 1000\,s post-mapping limit. For ePMapping this includes optimizer preparation, all rounds and final rebuild, excluding preprocessing, mapping and CEC. An outer 3600\,s guard covers setup and validation. Artifacts are archived.

Each configuration/design runs cold with an empty private pool, then warm with the same pool and limits. Cold may reload its own records after eviction or leave an incomplete pool on a timeout. Native caches restart; OS caches are not cleared. Both groups share the verified ABC batch. Times are $t_c/t_w$.

Completed outputs undergo CEC. At a deadline (TO), independent CEC and serialized area/delay audits validate the best saved implementation, or select the initial mapping. Snapshots consume the 1000\,s budget; post-run validation is outside it. ABC cannot export an interrupted result and retains its initial mapping on a timeout. Across all parts, cold/warm have \ColdCompleted{}/\WarmCompleted{} completions and \ColdMethodTimeouts{}/\WarmMethodTimeouts{} timeouts. Selected outputs meet their initial bounds. Of the timeouts, \ColdCheckpointSelections{}/\WarmCheckpointSelections{} select improved saved implementations. Of \ProfilePairs{} cold/warm pairs, \ProfileEqual{} agree in cost and \ProfileIdentical{} serialize identically. Differing values are stacked or separated by a slash, cold first. Independent audits include serialization buffers. Changes are $\delta A_j=100(A_j/A_{0,j}-1)$ and $\delta D_j=100(D_j/D_{0,j}-1)$; negative is better. Means include initial selections; ablations sum $\Delta A=\sum_j(A_{\rm arm,j}-A_{\rm full,j})$.

\subsection{Area Recovery after AIG Preprocessing}\label{sec:aspdac}

The preprocessed-input experiment measures area recovery after AIG preprocessing (\Cref{tab:main}). Both caps start from $N_e$. At one round, mean area changes are \OwnOneArea\%; mean delay changes are \OwnOneDelay\% (cold/warm). At cap 6, these are \OwnArea\% area and \OwnDelay\% delay, with \ColdMainCompleted{}/\WarmMainCompleted{} completed runs. Area decreases on \OwnAreaWins{} of sixteen; every selected result has CEC evidence and meets its initial bounds. Snapshots preserve partial progress after preprocessing. Independent cap-1/cap-6 runs can follow different trajectories under care budgets and deadlines.

For \texttt{arbiter}, area changes from $557.72$ to $\ArbiterArea\,\mu m^2$ and delay from $999.87$ to $\ArbiterDelay$\,ps. The other fifteen average \mbox{\OwnOtherArea\%} area and \mbox{\OwnOtherDelay\%} delay changes. Chain transactions occur in \ChainAreaWins{} of the completed designs. 



The ASPDAC'25 columns reconstruct final costs from its reported starts and rounded changes~\cite{aspdac25}; that implementation is not re-run. Our final area is smaller on \MainWins{} designs, with mean area and delay differences of \mbox{\MainAreaAll\%} and \mbox{\MainDelayAll\%}. These approximate end-to-end comparisons use different starts and cannot isolate post-mapping gains. Our delay bound refers to our input.

\subsection{Common-Input Comparison with ABC}\label{sec:mfs3}
The direct-input experiment uses the same mapping for both tools (\Cref{tab:mfs3}); ours uses one round. ABC's \texttt{mfs3 -a} does not enforce our initial-delay bound, so we distinguish raw and externally guarded outputs.

ABC completes \ABCCompleted{} designs and times out on \ABCTimeouts{}. On the ABC-completed subset, its mean area/delay changes are \mbox{\ABCArea\%}/\mbox{\ABCDelay\%}, versus \mbox{\PairedOwnArea\%}/\mbox{\PairedOwnDelay\%} for our verified selections. ePMapping completes \ColdMfsCompleted{}/\WarmMfsCompleted{} designs cold/warm. Its all-sixteen selected-output means are \mbox{\MfsOwnArea\%}/\mbox{\MfsOwnDelay\%}, within the initial delay bounds.

Of the completed ABC outputs, \ABCViolations{} violate the delay bound and \ABCFeasible{} satisfy both initial area and delay bounds. An external policy retains a completed result only if both hold, otherwise returning its initial mapping. Its paired means are \mbox{\ABCGuardArea\%} area and \mbox{\ABCGuardDelay\%} delay. This external selection leaves ABC unchanged; means describe the paired subset.

\Cref{tab:mfs3} shows the common mapping before ABC import. ABC may normalize output buffers, so its percentages use the parsed start; ours use the displayed costs. Times exclude CEC. ABC includes import and writing; ours includes optimizer preparation through rebuild. 

\subsection{Operator Ablations and Efficiency}\label{sec:oracles}
\AblationResultTable


All sixteen arms use identical $N_e$ inputs, bounds and round caps (\Cref{tab:ablation}); L1 remains enabled. noSYN disables all L3, including cofactors; noEGG disables L4. noDB removes lookup and database-derived mapper entries; noCARE supplies full care. noCHAIN skips composition; noSHARE clears chain memo tables between roots. noNARROW removes narrow initialization with unchanged extraction caps; noPHASE disables all physical phase reuse.

noSHANNON retains only Davio; noDAVIO retains only Shannon. noREUSE removes existing-INV fanout lookup but retains backward INV-input cancellation. noPRICE uses intrinsic area and timing filters while preserving physical reuse. Polarity-aware matching remains enabled; configurations are logged. The three construction controls in \Cref{fig:ablation} retain the candidate families: cache64K lowers each native cache limit to 65536; noAIGREUSE disables pre-mapping exact-AIG reuse; noLIBSHARE copies the cell library into each transient mapped network.

Removing Shannon and Davio gives summed area differences of \DeltaShannon{} and \DeltaDavio\,$\mu m^2$; removing forward reuse and live pricing gives \DeltaReuse{} and \DeltaPrice\,$\mu m^2$. \UsbAblationExample{} Full-ablation and main runs are independent measurements: deadlines and care budgets can change trajectories. \DeadlineExample{}

Removing the database improves \DBBetterCount{} designs; most aggregate gain comes from interrupted \texttt{sqrt} runs. noDB changes inner mapping too. Chain and sharing effects must not be added.

Removing L4 worsens area on \EggWorseCount{} designs; its cold/warm time ratios are \ColdNoEggRatio{}/\WarmNoEggRatio{}. L4 takes \TimeEggShare\%/\WarmTimeEggShare\% of recorded phase time in \ColdPhaseRuns{}/\WarmPhaseRuns{} completed full runs (CEC and other phases excluded). Optimizer totals at cap 6 are \ColdMainSeconds\,s/\WarmMainSeconds\,s. Persistent pool reuse avoids saturation and extraction on hits; native caches restart in each process. The cache64K time/full ratios are \CacheCapacityColdRatio{}/\CacheCapacityWarmRatio{}, versus \AigReuseColdRatio{}/\AigReuseWarmRatio{} for noAIGREUSE and \LibraryShareColdRatio{}/\LibraryShareWarmRatio{} for noLIBSHARE. \BuildCaseExample{} Avoided mapper calls need not reduce elapsed time: exact keys also cost time. These controls add no transformations; memory capacity and deadlines affect their results. \CacheDeadlineExample{} Ablations and shared-host scheduling affect elapsed time.

